% !TeX root = ../papers/nc-p3.tex
\section{Bell曲面の固有値衝突における全正規形・周期・成分交差}
\label{sec:split-repeated-surface}
\begin{theorem}\label{thm:split-repeated-surface}
可逆分解窓の曲面スペクトルで平方固有値に重複があり、直線スペクトルには属さないとする。
$t\ne0$, $t^4\ne1$, $f=(t^2-1)/(2t)$ として、以下の三行が全ての無印添字保持同型を与える。
三行は互いに交わらない。
\begin{enumerate}
\item 半単純・反転なしでは $u=f^2\in k\setminus\{0,-1\}$ が完全不変量である。
最小周期は $u=1$ で1、他で2である。
\item 半単純・反転ありでは $(u,v)=(f^2,fc)$, $c\ne0$ を取り、
$(u,v)\mapsto(u,u/v),(u^{-1},v/u)$ で割る。最小周期は $v^2=-u$ で1、他で2である。
\item 非半単純では $v=\pm f$, $v\ne0$, $v^2\ne-1$ 自身が完全不変量である。
全点の最小周期は2で、閉軌道は第二行の $(v^2,v)$ である。
\end{enumerate}
全点が標準Koszul左右AS正則である。十五次元成分に入るのは第一・第三行の全体と、
第二行の $v^2=u$ だけである。
\end{theorem}
\begin{proof}[全補完]
曲面の平方固有値は $1,r,s,rs$、$r,s\ne0,1$ なので、衝突は $r=s$ または $rs=1$ である。
共通倍率・置換・偶符号でBell四数を $(1,-t^2,t,t)$ に取れる。
実際、重なる二つの平方と残りの二平方の積の関係からこの規格化を得て、
行列式の符号は曲面式が決める。残り二平方の衝突は直線上の $t^2=-1$ である。
$P$ の通常二進成分は $0000,0011,0101$ と補語に順に $(1-t^2)/2,(1+t^2)/2,t$。
非半単純ならさらに $E_{0110}$ を加える。第一切断の奇ビット対平面で一鎖を作り、
残りの奇数長の鎖で向きを調整できるのでこれで十分である。

$P$ の安定群は、半単純なら $T(c,1/c)$ と全反転 $F$ の生成群、非半単純なら
そのうちの $T(\pm1,\pm1)$ だけである。実際、唯一の重固有平面とその二本の可分解線は
左右で単項作用を強制する。残りの二固有線は $00\pm11$ でそれぞれ保たれる。
従って正規化した対角作用は $T(c,1/c)$、反転は全ビット同時だけであり、付加語があれば
その係数と位置の保存から後者が消え $c^2=1$ となる。
$D_0=(Z,1,Z,1)$ として、半単純の全第二対応は $T(c,1/c)F^eD_0$ である。
$D_0$ が移送になるのは定理\ref{thm:split-surface-component}の $h=1$ の表示による。
対角基底は $e=0$ で $c$ を自由に変え、$e=1$ では自明に作用し、$F$ は $c$ を逆にする。
付加語があるときも重固有平面から全局所行列は単項で、両逆テンソルの奇ビット対ブロックは
三角になる。その唯一の完全マッチングを保つため半単純部も移送する。
付加項は $I_0,I_1$ に順に $-t^{-2}E_{1001},-4(t^2-1)^{-2}E_{1001}$ として現れる。
両切断で比較すると $e=1$, $c^2=1$ が必要十分である。
この表示を $T(z,z^{-1})$ の比を用いて付加項だけ消せば同じ $t,c$ の半単純部へ流入する。
半単純部の閉性は、各ブロック内の補語に同じ正重み、各ブロックに等しい全重みを与えれば
一様な周辺を持ち、一文字だけ異なる非零語もないことから従う。
\end{proof}

\begin{proof}[無印同型]
反転なしは対角移送に対する固定群が有限なので全安定群の階数1であり、指定分解は交換を除き一意である。
逆数Bell四数に移る交換は $t\mapsto1/t$ となる。重固有値の規格化と二単純固有線の交換から
許される $t$ は $\pm t^{\pm1}$ であり、これらは全て同型になる。
実際、元の四数を共通倍率、置換、偶符号で移せば、移送側は階数1の唯一の半単純補完である。
よって不変量は $u$ である。一移動と分解交換では旧 $P$ を一巡回するため、
新しいBell四数は $((1-t^2)/2,(1-t^2)/2,(t+1)^2/2,-(t-1)^2/2)$ である。
従って $u$ は $1/u$ に移り第一行の周期を得る。

反転ありでは、位置2,3でビットを反転する。非零語は24置換語に
$0202,2020,1313,3131$ を加えたもの、Jordanの場合はさらに $0303,3030$ を加えたものである。
半単純では全位置共通の対角トーラス $l_0+l_2=l_1+l_3$ が極大で各位置に異なる四指標を持つため、
全同型を単項同型で比較できる。非半単純では階数1で分解が符号まで一意、
重固有平面の二本の可分解線が各局所二文字の座標線を決めるため、交換の場合を含めやはり単項で足りる。
台で隣接対として現れないのは等しい二文字だけなので、位置別の許される置換は四位置共通である。
従って第二行ではマッチング $02,13$ を保つ八置換、第三行ではさらに辺 $03$ を保つ
恒等と $(03)(12)$ の二つに限られる（四頂点における二端の次数も保存する）。

対角同値を具体的に計算する。$w_q$ を語 $q$ の係数として、二つの不変量は
\[
 u=\frac{w_{0213}w_{1302}}{w_{0312}w_{1203}},\qquad
 v=-\frac{w_{2031}w_{0132}w_{1230}}{w_{0231}w_{2130}w_{1032}} .
\]
六ブロックの式でこれらは $f^2,fc$ となる。
\path{computations/verify_split_repeated_surface.py} は位置別文字の28列の整数重み行列の核を取り、
各核指標の値を $t,t-1,t+1,t^2+1,c$ の指数ベクトルに写す。
像の格子の基底は $(-2,2,2,0,0),(-2,0,0,2,0),(-1,1,1,0,1)$ である。
従って指標値は定数を除き $u,1+u,v$ の冪積で、同じ $u,v$ で元の二表示は対角同値である。
八置換を行うと対角同値の範囲で $t$ は $t,(t+1)/(t-1)$ に、$c$ は $c,c^{-1}$ に独立に移る。
上の検算は各比が全整数核上で1であることを符号も含めて確認する。
これが表示した二つの生成変換である。一移動を合わせるとさらに $c$ が $-c$ になる。
従って無印では $v\mapsto-v$ であり、$u=-1$ は除外されているから第二行の周期条件を得る。

Jordanの場合、二置換による第二行の比較では $u,v$ ともに保たれる（$c=\pm1$）。
付加二語の相対倍率について、前節の辺 $03$ の六係数比はやはり $-1$ である。
よって $u,v$ が等しければ、マッチングの対角トーラスの残りの非自明指標で
二項共通の倍率も一致させられる。これで第三行を得る。一移動同型は閉極限の周期から存在せず、
二移動では付加語を含む全整数核上でも係数比が1となる。
最後に第二行の別分解の可能性を確かめる。その極大トーラスの四指標から、
二重みずつの分解は $01\mid23$ と $03\mid12$ だけである。
後者の交代部分の $00,11$ ブロックは反対角で二係数の積 $f^2$、
$01,10$ ブロックの二つの平方固有値は $((t^2+1)/(2t)\pm1)^2$ である。
従ってこちらの印でも直線外の同じ重複型である。第一・第三行では印が符号まで一意で、
逆転置もこの重複型とJordan分割を保つ。これで直線部や単純平方スペクトルとの重複も起きない。
\end{proof}

\begin{proof}[十五次元成分との全交差]
反転なしと $c=\pm1$ の半単純窓は曲面の四族に $u=v=t$ を代入して得られる。
Jordanの二表示も同じ成分に入る。実際、曲面一般の四数 $(1,-xy,x,y)$ において
$(x-y)/2=z^2$、$x,y\to t$ とし、能動側に $T(z,z^{-1})$ を作用させる。
$0110$ の係数だけが1として残り、補語は零になる。$F$ を伴う二表示の四対応は
この対角変更で不変なので所定のJordan窓になる。

逆向きに半単純反転窓では、両色の純軸で純二文字関係を $r_i$ としたとき、
$Jr_i^t$ の四移動の積は射影的に $\operatorname{diag}(1,c^{-2})$ である。
これは $T(c,1/c)FD_0$ の四対を純ブロックで順に取った式である。
従って十五次元成分に入るなら定理\ref{thm:pi-periodic-point-paths}により、
$c^2$ の位数は $1,2,3,4$ である。$c^2\ne1$ の候補の最小多項式は
$c^2+1,c^2+c+1,c^2-c+1,c^4+1$ に限る。
各剰余体上で隣接二切断の接方程式を取り、$t(t^4-1)\ne0$ の全域で
アフィン接次元は65である。具体的に各辺の六降順対をピボットに取ると
残り十座標から核基底が得られる。ピボットの可逆性は純二式が交換式、混合四式が
二方向の間の同型であることから全域で成立する。
二切断のそれぞれについて両核基底の100積を作り、トーラスの二指標別の25ブロックに分ける。
上の検算は表示中の分母が指定開集合で単元であることを調べ、
一変数多項式環上のSmith形の非零不変因子が全てこの開集合で単元となること、
零方向の次元和が65であることを確認する。十五次元成分のアフィン次元76を下回るため、
これらの候補は例外 $t$ を含め全て除かれる。
\end{proof}
